% Nouveau EGA V (cote 25): archive cote 25, site page 142 (PDF page 143).
% typescript. Draft compared: none.
% AI-transcribed by Claude Opus 5.5 (claude-opus-5-5, adaptive thinking, effort high), Message Batches API, batch ega25b (msgbatch_01A7n5YuXubUPNK1zsepi8RG), 2026-11-13; not yet checked by eye.
\scanpage{25/142}
\mspage{\hand{3}}

préschéma x sur k, \xout{enfin la dernière assertion résulte de la première}
\xout{précédente et de l'isomorphisme $\underline{m}_x$ analogue $gr^2(P^2_{X/k})(x)$ Sy}

Notons que $gr^1(\underline{P}^1_{X/Y}) \simeq \underline{\Omega}^1_{X/Y}$ , par suite la restriction
$d^1\phi\,|V(\phi)$ peut être interprété comme une section de $\underline{\Omega}^1_{X/Y} \otimes \underline{O}_{V(\phi)}$ ,
qui n'est autre d'ailleurs que la restriction de $d_{X/Y}\,\phi$ à $V(\phi)$ . Nous
pouvons alors considérer le \uline{préschéma des zéros de cette section} , que
nous noterons $V(\phi)^{sing}$ , et dont l'ensemble sous-jacent n'est autre
que l'ensemble des zéros de $\phi$ singuliers rel à Y en vertu de prop.5 .
( NB si \xout{\unsure{$\phi$}} $\mu$ est une section d'un Module localement libre \inserted[typed]{de type fini} $\underline{E}$ sur un
\margin{$\|$}\inserted[typed]{préschéma} \xout{espace annelé} X, on définit de façon évidente le sous-préschéma d es
zéros de $\mu$ , par exemple comme défini par l'Idéal image de $\check{\underline{E}} \longrightarrow \underline{O}_X$
transposé de $\mu$ ; lorsque $\underline{E} = \underline{O}_X{}^n$ et $\mu = \xout{\unsure{(\mu_i)}}\,(\mu_1,..,\mu_n)$, alors cet
idéal n'est autre que $\sum \mu_i\,\underline{O}_X$ ,qui définit $V(\mu_1,..,\mu_n)$ ). Prenant
maintenant la restriction $d^2\phi\,|V(\phi)^{sing}$ , et notant que
$gr^2(P^2_{X/Y}) \simeq Sym\xout{m}^2(\underline{\Omega}^1_{X/Y})$, on trouve une section canonique $M(\phi)$ de
$Sym^2(\underline{\Omega}^{\,1}_{X/\unsure{Y|X}}) \otimes \underline{O}_{V^{sing}}$ . On vérifie immédiatement que (prenant des
\margin{$\|$}point de X à val\xout{eu}rs dans des corps ...) cette section est précisément
celle qui détermine les formes quadratiques envisagées dans déf 1
(dans le cas \xout{\illegible{1}} d'un $X_k$ déduit de X/S par Spec(k)$\rightarrow$ S). On
en déduit une description d e l'ensemble $V(\phi)^{sup\ sing}$ en termes de
cette section, de la façon suivante: interprétant $M(\phi)$ comme définissan
% NOTE: the end of this line is cut by the right edge of the scan (presumably "définissant" / "définissant un")
homomorphisme
\[
M(\phi)\hand{$'$} : \ \hand{$\mathcal{T}$}_{X/Y} \otimes_{\underline{O}_X} \underline{O}_{V^{sing}} \longrightarrow \underline{\Omega}^1_{X/Y} \otimes_{\underline{O}_X} \underline{O}_{V^{sing}} \quad ,
\]
il faut prendre l'ensemble des points en lesquels cet homomorphisme
n'est pas un isomorphisme. Cela montre en particulier que $V(\phi)^{sup\ sing}$
est un ensemble fermé. On peut \deleted{le} préciser \inserted{ce dernier}\deleted{, en} introduisant
\[
D(\phi) = \text{dét}\ M(\phi) \ \in \ \Gamma\,(\underline{\Omega}^d_{X/Y}{}^{\otimes 2} \otimes \underline{O}_{V^{sing}}\ ) \quad ,
\]

\mspage{\hand{142}}
